\begin{lemma}
For every fixed \(\delta>0\) and \(\varepsilon>0\), there is a threshold
\(n_0\) such that the following holds for every \(n\ge n_0\).  Let
\[
m=\noderef{TwoBiteNaturalReducedVertexCount}(n)
=\left\lfloor\frac{n}{\log^2 n}\right\rfloor,
\qquad
k=\noderef{TwoBiteNaturalIndependenceScale}(1+\varepsilon,n)
  =\left\lceil(1+\varepsilon)\sqrt{n\log n}\right\rceil .
\]
For a fixed set \(I\subseteq V(G)\) with \(|I|=k\), and integers
\(\ell_R,\ell_B\) with \(\ell_R\ell_B\ge k\), put
\[
x_R=\frac{\ell_R}{k},\qquad x_B=\frac{\ell_B}{k}.
\]
Under the two-bite law, the probability of the event
\[
\noderef{TwoBiteProjectionSizeEvent}(I,k,\ell_R,\ell_B)
\]
is at most
\[
\exp\left(\left(\delta-\frac{2-x_R-x_B}{2}\right)k\log n\right)
\]
for all sufficiently large \(n\).
\end{lemma}
\begin{proof}
Fix \(\delta>0\) and \(\varepsilon>0\).  We choose \(n_0\) large enough for all
large-\(n\) comparisons below.  In particular, for every \(n\ge n_0\) the real
quantities
\[
M_{\mathbb R}(n)=\noderef{TwoBiteReducedVertexCount}(n)=\frac{n}{\log^2 n},
\qquad
K_{\mathbb R}(n)=\noderef{TwoBiteIndependenceScale}(1+\varepsilon,n)
  =(1+\varepsilon)\sqrt{n\log n}
\]
are positive, and \(\noderef{NaturalParameterApproximation}\) gives
\[
m\le M_{\mathbb R}(n)<m+1,\qquad
K_{\mathbb R}(n)\le k<K_{\mathbb R}(n)+1.
\]
Increasing \(n_0\) further, we may use the consequences
\[
m\ge \frac12\,\frac{n}{\log^2 n},\qquad
k\le 2(1+\varepsilon)\sqrt{n\log n},\qquad
1\le k\le n\le m^2.
\tag{2}
\]
The first two inequalities follow from \(M_{\mathbb R}(n)<m+1\) and
\(k<K_{\mathbb R}(n)+1\) after requiring
\(M_{\mathbb R}(n)\ge2\) and \(K_{\mathbb R}(n)\ge1\).  The inequalities
\(1\le k\le n\le m^2\) then follow by increasing \(n_0\) so that
\[
2(1+\varepsilon)\sqrt{n\log n}\le n
\quad\text{and}\quad
n\le\left(\frac{n}{2\log^2 n}\right)^2.
\]
The last inequality ensures that injections
\(\noderef{TwoBiteEmbedding}:V(G)\hookrightarrow V_R\times V_B\) exist and that
the \(k\)-point image of \(I\) can be compared with all \(k\)-subsets of
\(V_R\times V_B\).

We first reduce \(\noderef{TwoBiteEventProbability}\) to a counting problem for
uniform injections.  A sample
\(\omega\in\noderef{TwoBiteSample}(n,m,p)\) consists of a red graph \(G_R\), a
blue graph \(G_B\), and an injection \(\pi\).  By
\(\noderef{TwoBiteSampleWeight}\), its weight is
\[
\noderef{GnpGraphWeight}(p,G_R)\,
\noderef{GnpGraphWeight}(p,G_B)\,
\noderef{UniformInjectionWeight}(n,m).
\]
The event \(\noderef{TwoBiteProjectionSizeEvent}(I,k,\ell_R,\ell_B)\) depends
only on the injection coordinate, because
\(\noderef{TwoBiteRedProjection}\) and \(\noderef{TwoBiteBlueProjection}\) are
the two coordinate projections of \(\pi=\noderef{TwoBiteEmbedding}(\omega)\).
For a fixed injection \(\pi\), summing \(\noderef{GnpGraphWeight}\) over all red
graphs gives \(1\): each unordered edge contributes the two alternatives
\(p\) and \(1-p\), whose sum is \(1\), and the finite product over edges
therefore sums to \(1\).  The same argument applies to the blue graph.  Since
\(n\le m^2\), \(\noderef{UniformInjectionWeight}(n,m)\) is the reciprocal of
the number \(N_{\mathrm{inj}}\) of injections \(V(G)\hookrightarrow
V_R\times V_B\).  Thus
\[
\noderef{TwoBiteEventProbability}(n,m,p)(\mathcal E)
=\frac{\#\{\pi:\pi\text{ is an injection and }\mathcal E(\pi)\}}
       {N_{\mathrm{inj}}},
\]
where \(\mathcal E\) is the projection-size event for \(I\).  This also proves
that the left side is at most \(1\).

For the fixed \(k\)-set \(I\), the image set \(\pi(I)\) is uniformly
distributed over all \(k\)-subsets of \(V_R\times V_B\).  Indeed, if
\(S\subseteq V_R\times V_B\) has \(|S|=k\), then the number of injections
\(\pi\) with \(\pi(I)=S\) is
\[
k!\,(m^2-k)(m^2-k-1)\cdots(m^2-n+1),
\]
where the falling product counts the completion on \(V(G)\setminus I\).  This
number is independent of \(S\).  Consequently all \(\binom{m^2}{k}\)
\(k\)-subsets are equally likely as \(\pi(I)\).

We now count image sets with the requested projection sizes.  If
\(\ell_R>m\) or \(\ell_B>m\), then no red or blue projection set of that size
exists, so the event is empty.  If \(\ell_R>k\) or \(\ell_B>k\), then the event
is also empty, because a \(k\)-element image set has at most \(k\) red
coordinates and at most \(k\) blue coordinates.  These empty cases satisfy the
desired inequality trivially.  We may therefore assume
\[
1\le \ell_R,\ell_B\le m,\qquad
\ell_R,\ell_B\le k.
\]
The positivity follows from \(1\le k\) and the hypothesis \(\ell_R\ell_B\ge k\).

Choose first the red projection set \(R\subseteq V_R\) and the blue projection
set \(B\subseteq V_B\), with
\[
|R|=\ell_R,\qquad |B|=\ell_B.
\]
There are \(\binom m{\ell_R}\binom m{\ell_B}\) choices.  Once \(R\) and \(B\)
are fixed, every image set with these exact projections is a \(k\)-subset of
\(R\times B\).  Replacing "exactly \(R,B\)" by the weaker condition
"\(S\subseteq R\times B\)" overcounts, and gives at most
\(\binom{\ell_R\ell_B}{k}\) image sets for this pair \(R,B\).  Since
\(\ell_R\ell_B\ge k\), this binomial coefficient is defined in the nonempty
range.  Hence
\[
\Pr(\mathcal E)
\le
\binom m{\ell_R}\binom m{\ell_B}
\frac{\binom{\ell_R\ell_B}{k}}{\binom{m^2}{k}}.
\tag{3}
\]
If \(\Pr(\mathcal E)=0\), then the desired inequality is immediate because the
right-hand side is a positive exponential.  Thus, in the remaining argument we
may assume \(\Pr(\mathcal E)>0\).

We split into two cases based on the defect \(s=(2-x_R-x_B)k\).  If \(s\le2\delta k\), then
\[
\left(\delta-\frac{2-x_R-x_B}{2}\right)k\log n
=\left(\delta k-\frac{s}{2}\right)\log n\ge0
\]
for \(n\ge n_0\).  The exponential upper bound is then at least \(1\), while
the reduction to uniform injections already showed \(\Pr(\mathcal E)\le1\).
This proves the claim in the boundary regime.

Assume now that \(s>2\delta k\).  The required bound then follows from
\(\noderef{ProjectionSizeAnalyticBinomialBound}\), which bounds the binomial quotient in (3)
by the same exponential when the defect \(s=(2-x_R-x_B)k\) is strictly greater than \(2\delta k\).
\end{proof}
